\documentclass[11pt,a4paper]{article}

\usepackage[utf8]{inputenc}

\usepackage[T1]{fontenc}

\usepackage[polish]{babel}

\usepackage{lmodern}

\usepackage{geometry}

\geometry{a4paper,margin=27mm}

\usepackage{microtype}

\usepackage[hidelinks]{hyperref}

% Portfolio: MSC-AR-P22; source: prospective ARC-XR research scope.

\title{Neuroadaptive Social XR z odrębną personalizacją uczestników współpracy}

\author{Mateusz Pomianek}

\date{Wersja 0.3.0, 21 listopada 2026}

\begin{document}

\maketitle

\begin{abstract}

Przedmiotem opracowania jest Neuroadaptive Social XR z odrębną personalizacją uczestników współpracy. Czy indywidualna adaptacja w Social XR poprawia koordynację względem wspólnej strategii dla całej grupy? Badanie zestawi wspólną strategię adaptacji dla zespołu z indywidualnymi profilami dostosowanymi do poszczególnych uczestników. Głównym punktem końcowym jest liczba błędów koordynacji i czas wykonania wspólnego celu. Opis ma charakter wstępny i nie przedstawia niezweryfikowanych wyników.

\end{abstract}

\section{Wprowadzenie}

Rozwój adaptacyjnych systemów XR stwarza możliwość badania interakcji człowieka z informacją przestrzenną, ale samo zastosowanie AR, VR lub modeli sztucznej inteligencji nie przesądza o przewadze poznawczej czy ergonomicznej. W rozpatrywanym zagadnieniu, Neuroadaptive Social XR z odrębną personalizacją uczestników współpracy, kluczowe jest rozdzielenie właściwości systemu od zmienności zachowania użytkownika. Porównanie należy zaprojektować tak, aby mierzyć jednoznaczne zjawisko, a nie tylko ogólną ocenę technologii.

\section{Cel pracy i pytanie badawcze}

Czy indywidualna adaptacja w Social XR poprawia koordynację względem wspólnej strategii dla całej grupy? Projekt wymaga krytycznego przeglądu aktualnego stanu badań oraz zdefiniowania nowości w stosunku do innych tematów portfolio XR.

\section{Zakres i odrębny wkład naukowy}
Indywidualizacja adaptacji w Social XR wobec wspólnej polityki grupowej. Wynik jednostkowy nie może ukrywać kosztu obciążenia członka zespołu. Odrębność dotyczy zespołowego błędu koordynacyjnego i kosztu nierówności indywidualnych korzyści, nie ogólnej jakości Social VR.
\section{Hipotezy badawcze}
\paragraph{H1.} Indywidualne profile adaptacyjne zmniejszają liczbę konfliktów koordynacji względem jednolitej polityki zespołu.
\paragraph{H2.} Korzyść pozostaje po kontroli czasu wykonania i jednakowej dostępności informacji dla każdego uczestnika.
\paragraph{H3.} Polityka poprawia wynik zespołu bez pogorszenia najsłabszego indywidualnego wyniku powyżej z góry określonego marginesu.
\section{Projekt badania i plan eksperymentów}
\subsection{A. Operacjonalizacja i kwalifikacja danych}
Indywidualizacja adaptacji w Social XR wobec wspólnej polityki grupowej. Wynik jednostkowy nie może ukrywać kosztu obciążenia członka zespołu. Jednostka wnioskowania: zespół współpracujących uczestników i niezależna sesja wspólnego zadania. Przed badaniem należy określić źródła, wersje urządzeń, czas akwizycji, etykiety i niezależną referencję. Nie liczyć zależnych okien i klatek jako niezależnych prób.
\subsection{B. Eksperyment porównawczy}
Losować kolejność zespołowych warunków bez zmiany składu zespołu w ramach sesji, zachować te same role i współdzielone cele. Wynik pierwszorzędowy: różnica liczby błędów koordynacji na ukończony wspólny cel: indywidualna adaptacja vs jedna adaptacja grupowa. Zamrozić losowanie warunków, progi, mianownik punktu końcowego i kryteria wyłączeń bez dostępu do oceny końcowej.
\subsection{C. Kontrole i analiza odporności}
Rejestrować przydział informacji, konflikt między uczestnikami i odsetek opóźnionych komunikatów; nie ujawniać prywatnych profili członkom zespołu. Ustalić kryteria sprawiedliwości dostępu dla mniej sprawnego użytkownika; osobno mierzyć obciążenie i niezadowolenie. Utrata sygnału, abstencja, przerwanie zadania i dane brakujące są odrębnymi kategoriami wyniku.
\subsection{D. Walidacja na niezależnym materiale}
Walidacja musi zachować podział po uczestniku, sesji, scenie lub zadaniu. Model liczby błędów z klastrowaniem zespołowym, analiza interakcji uczestnik × polityka i raport wyników najgorszego członka. Raportować wyniki negatywne i rozbieżności z założeniami.
\section{Estymand, analiza statystyczna i decyzje}
różnica liczby błędów koordynacji na ukończony wspólny cel: indywidualna adaptacja vs jedna adaptacja grupowa. Jednostka: zespół współpracujących uczestników i niezależna sesja wspólnego zadania. Model liczby błędów z klastrowaniem zespołowym, analiza interakcji uczestnik × polityka i raport wyników najgorszego członka. H1 jest planowanym kontrastem głównym; H2 i H3 pozostają pomocnicze do zatwierdzenia hierarchii testów oraz kontroli wielokrotności. Wartości marginesów, mocy, minimalnego efektu istotnego praktycznie i liczebności mają wynikać z pilotażu, a nie z post hoc wyboru wyniku.
\section{Stan badań i granica nowości}
\begin{itemize}
\item Multiuser VR research protocols exist, so interaction-format effects must be isolated under controlled social and scene conditions. \cite{xrprior1}.
\item XR personalization methods exist; a claim needs a defined multi-session or transfer contrast. \cite{xrprior2}.
\end{itemize}
Odrębność dotyczy zespołowego błędu koordynacyjnego i kosztu nierówności indywidualnych korzyści, nie ogólnej jakości Social VR. Odrębność na tle opublikowanych metod i pozostałych artykułów wymaga osobnego audytu.
\section{Warunki bezpieczeństwa, sprzętu i zgodności}
To prospektywna koncepcja, nie zgoda na eksperyment, prerejestracja lub wynik empiryczny. Praca z uczestnikami, interwencje stresowe i pomiar biometryczny wymagają przeglądu etycznego, świadomej zgody i ochrony danych. Quest 3 w wersji standardowej nie posiada eye trackingu, a Tobii Eye Tracker 5 jest urządzeniem ekranowym; nie zastępować jednego stanowiska drugim.
\section{Materiał i metody}

Badanie zestawi wspólną strategię adaptacji dla zespołu z indywidualnymi profilami dostosowanymi do poszczególnych uczestników. Wymagane jest zdefiniowanie jednostki analizy, zasad doboru przypadków, wersji scenariusza i źródeł pomiaru. Rozkład warunków oraz reguły wykluczenia prób należy określić przed pomiarem. W badaniach z udziałem ludzi niezbędna jest uprzednia właściwa ocena etyczna i świadoma zgoda. W testach wyłącznie symulacyjnych należy oddzielić prawdę symulacyjną od obserwacji empirycznej.

\subsection{Wynik główny i kryteria jakości}

Liczba błędów koordynacji i czas wykonania wspólnego celu. Pomiar powinien zachowywać jawne jednostki, znaczniki czasu i reguły oceny poprawności. Przy modelach predykcyjnych stosować rozłączne partycje oceny; przy eksperymentach interwencyjnych oddzielić skuteczność predykcji od przyczynowego wpływu interfejsu.

\section{Plan analizy}

Model uwzględniający zagnieżdżenie uczestników w zespołach oraz asymetrię obciążenia. W wynikach należy raportować wielkość efektu lub błąd estymacji z przedziałem niepewności, liczbą obserwacji i odsetkiem odrzuconych prób. Analizy dodatkowe wymagają jasnego oznaczenia jako eksploracyjne. Nie wybierać progów decyzyjnych ani modelu na podstawie obserwacji ze zbioru testowego.

\section{Wyniki}

W obecnej wersji nie raportuje się wyników pomiarów, tabel statystycznych ani parametrów wyznaczonych z nieudokumentowanych danych. To kierunek badawczy, a nie zatwierdzony protokół badania ani zgłoszony artykuł.

\section{Dyskusja}

Interpretacja przyszłych wyników musi uwzględnić alternatywne mechanizmy i zależność wniosków od jakości danych, warunków zadania oraz ograniczeń sprzętu. Wartość potencjalnego rozwiązania zależy nie tylko od trafności modelu, lecz także od jego opóźnienia, stabilności i wpływu na zachowanie człowieka. Wyniki testów syntetycznych nie mogą być automatycznie uogólniane na populację uczestników.

\section{Ograniczenia badania}

Poprawa globalnego czasu nie może ukrywać pogorszenia dostępności lub obciążenia części uczestników. Granice nowości i trafności ekologicznej wymagają niezależnej weryfikacji. Żadna część szkicu nie dowodzi przeprowadzenia zbierania danych, uzyskania zgód etycznych ani gotowości manuskryptu do zgłoszenia.

\section{Wnioski}

Przedstawiono problem, kontrast i sposób jego planowanej oceny. Rozstrzygnięcie postawionego pytania wymaga przygotowania protokołu, zweryfikowanych źródeł oraz poprawnie udokumentowanych wyników. Bez tych materiałów nie można sformułować konkluzji empirycznych.

\begin{thebibliography}{9}
\bibitem{xrprior1} Creating a social virtual reality application for psychological research: A tutorial (2025), Behavior Research Methods. \url{https://doi.org/10.3758/s13428-025-02693-4}.
\bibitem{xrprior2} Adapting to the User: A Systematic Review of Personalized Interaction in VR (2025), arXiv:2510.13123. \url{https://arxiv.org/abs/2510.13123}.
\end{thebibliography}
\end{document}

